\chapter{Training Loop and Python Integration}\label{ch:training-and-python}

\section{Overview}

Two consumers drive the trainable pathway (Chapter~\ref{ch:tile-renderer}):
the native Rust training loop in \code{rmesh-train}, and the PyTorch bridge
in \code{rmesh-python}. Both share the same principle: record forward,
loss, backward, and the optimizer into as few command encoders as possible
and keep data on the GPU --- the native loop's only per-step readback is the
scalar loss.

\section{The native loop (\code{rmesh-train})}

\code{train()} builds every pipeline, buffer, and bind group once (including
A/B variants for the radix-sort ping-pong), then per view records a
\emph{single} command encoder:

\begin{enumerate}
\item clear gradient, loss, tile-range, and image buffers;
\item \code{project_compute} --- SH/color evaluation, culling, per-tet tile
  counts, compact visible ids;
\item tile-pair generation and 64-bit radix sort, then tile ranges
  (Chapter~\ref{ch:gpu-sorting});
\item forward rasterization --- the native loop uses the
  software-rasterizer baseline (tile size 12,
  \code{RasterizeComputePipeline}); the interval variant is exercised
  through the Python bindings;
\item loss pass; backward pass; three Adam dispatches (one per parameter
  group).
\end{enumerate}

One \code{queue.submit} per view; the only \code{write_buffer}s are the
camera uniforms and per-group Adam uniforms. Cameras come from
\code{TrainingView} with the pinhole convention (x right, y down, z
forward) and intrinsics derived from the FOV.

\subsection{Loss (\code{loss_compute.wgsl})}

A $16\times16$ workgroup per pixel block compares the premultiplied rendered
RGB against ground truth. \code{loss_type = 0} gives L1
($\partial L/\partial c = \text{sign}(c - g)/n$), otherwise L2
($2(c - g)/n$); the result is written per pixel to \code{dl_d_image}
(alpha gradient zero). SSIM has plumbing (a \code{lambda_ssim} uniform and
\code{ssim_err} buffers consumed by \code{rmesh-error}) but is \emph{not}
computed by this shader --- only L1/L2 are.

\subsection{Adam (\code{adam_compute.wgsl})}

One dispatch per parameter group so each group gets its own learning rate:
vertices ($10^{-4}$), densities ($10^{-2}$), color gradients ($10^{-3}$);
$\beta_1 = 0.9$, $\beta_2 = 0.999$, $\epsilon = 10^{-8}$. The kernel is a
standard Adam step with NaN-guarded gradients and bias correction via
$\beta^t = e^{t \ln \beta}$, using the safe-math helpers throughout. Moment
buffers ($m, v$) live per group in \code{AdamState} /
\code{MaterialAdamState}.

\subsection{Per-tet error statistics (\code{rmesh-error})}

\code{error_tiled_compute.wgsl} re-traverses tiles exactly like the forward
rasterizer but accumulates per-tet statistics with atomics, weighted by
$T \cdot \tau$: total weight, weighted L1 and L1$^2$, weighted SSIM error,
weighted entry/exit points and centroid, a hit count, and a max-visibility
score. These drive densification decisions (which tets to split or prune).

\section{Python bindings (\code{rmesh-python})}

\code{maturin develop} builds the \code{_native} module; the
\code{rmesh_wgpu} package wraps it.

\subsection{\code{RMeshRenderer}}

A single \code{#[pyclass]} owns the entire wgpu state. Key methods:

\begin{itemize}
\item constructor: \code{(vertices, indices, base_colors,}
  \code{densities, color_grads, circumdata, width, height, aux_dim=0)};
\item forward variants: \code{forward} (hardware-raster path),
  \code{forward_tiled} (software rasterizer),
  \code{forward_interval_tiled} (interval path --- returns a dict with
  image, xyzd, distortion, and aux arrays), \code{forward_raytrace},
  \code{forward_error};
\item backward variants: \code{backward} and
  \code{backward_interval_tiled}, the latter taking \code{dl_d_image} and
  optional \code{dl_d_xyzd} / \code{dl_d_distortion} / \code{dl_d_aux}
  --- each returns a dict of flat
  \code{f32} gradient arrays (\code{d_vertices}, \code{d_densities},
  \code{d_base_colors}, \code{d_color_grads}, \code{d_vertex_normals},
  optionally \code{d_aux_data});
\item \code{train_step(...) -> f32} --- the all-GPU step (forward + loss +
  backward + Adam) with only the scalar loss read back;
\item \code{update_params} / \code{get_params} and setters for vertex
  normals, aux data, and \code{min_t}.
\end{itemize}

\code{backward_interval_tiled} internally re-runs tile generation and
sorting for the given view, then the two backward stages (per-tile
reverse traversal, then the per-tet chain back to parameters) before reading
back the gradient buffers.

\subsection{\code{autograd.py} and the tensor boundary}

\code{RMeshForward(torch.autograd.Function)}:
\begin{itemize}
\item \textbf{forward}: syncs the current PyTorch parameters into the GPU
  buffers via \code{update_params}, calls the chosen forward, and wraps the
  returned image back into a torch tensor on the caller's device;
\item \textbf{backward}: moves \code{grad_output} to CPU numpy, calls the
  matching native backward, and wraps each gradient array as a torch tensor
  aligned with the forward argument tuple.
\end{itemize}

The boundary is \emph{explicit host copies} in both directions
(\code{.detach().cpu().numpy()} up, \code{read_buffer_f32} $\to$ numpy
$\to$ \code{torch.from_numpy} down); there is no zero-copy
CUDA$\leftrightarrow$wgpu interop. \code{RMeshModule(nn.Module)} exposes
\code{vertices}, \code{base_colors}, \code{densities}, and
\code{color_grads} as \code{nn.Parameter}s so a standard
\code{torch.optim.Adam(model.parameters())} loop trains the mesh --- the
alternative to the fully-native \code{train_step}.

\section{Validation}

Finite-difference gradient checks live in
\code{crates/rmesh-render/tests/gradient_test.rs} and
\code{crates/rmesh-trainable/tests/kernel_tests.rs}; the analytic math they
check is derived symbolically in \code{scripts/derive_backward.py}
(see Chapter~\ref{ch:tile-renderer}). Renderer-vs-CPU-reference image comparisons
live in \code{crates/rmesh-render/tests/}.

